problem:  
Let \( X \) be a *Fano threefold* with smooth anticanonical divisor \( D \in |-K_X| \).  
For each positive integer \( k \), let  
\[
\pi_k : Y_k \to X
\]  
be the smooth cyclic cover branched of order \( k \) along \( D \), so that the canonical bundle satisfies \( K_{Y_k} = \pi_k^*(K_X + (k-1)/k \, D) \), hence for \( D \in |-K_X| \), we have \( K_{Y_k} \cong \mathcal{O}_{Y_k} \). Thus each \(Y_k\) is a smooth compact **Calabi–Yau** threefold.  

By Yau’s theorem, each \(Y_k\) admits a unique Ricci–flat Kähler metric \( g_k \) in any prescribed Kähler class \( [\pi_k^*\omega] \).  
Let \( D_k = \pi_k^{-1}(D) \) denote the ramification divisor in \( Y_k \), and note that locally near \( D_k \), the covering structure introduces a scale of branching that magnifies neighborhoods of \(D\) as \(k \to \infty\).

**Question:**  
After rescaling each metric \(g_k\) by an appropriate factor (for example, fixing unit volume or the mean curvature scale along \( D_k \)), do the sequences  
\[
(Y_k, g_k)
\]
converge, in a pointed Cheeger–Gromov or Gromov–Hausdorff sense, to a complete noncompact Ricci‑flat Kähler manifold  
\[
(Y_\infty, g_\infty)
\]
that is asymptotically conical (AC) with asymptotic cone determined by the Sasaki–Einstein link \( L \) associated to the anticanonical divisor \( D \)?  

Equivalently: does this **canonical Calabi–Yau cyclic‑cover sequence of Fano manifolds** give a natural compact‑to‑noncompact degeneration mechanism whose local limit near the branching reproduces an asymptotically conical Calabi–Yau geometry modeled on \( K_D^\times \) or its Reeb–Sasaki–Einstein cone structure?

If such a limit \( (Y_\infty, g_\infty) \) exists, is it unique up to holomorphic isometry, and can one detect the complex‑algebraic data of \( (D, -K_X|_D) \) directly from the asymptotic cone geometry?

---

Why is it a "good" problem:  
This version corrects the canonical‑bundle inconsistency and rests on a coherent, well‑defined algebraic setup: taking cyclic branched covers of **Fano** manifolds along *anticanonical* divisors produces genuine compact Calabi–Yau manifolds. It therefore provides a genuinely sound starting point for analytic inquiry.

Analytically, the problem investigates the **asymptotic metric behavior** of Ricci‑flat Kähler metrics under algebraically controlled *degeneration sequences*. It asks whether increasing the branching order along an anticanonical divisor drives a compact Calabi–Yau sequence toward a **noncompact Tian–Yau–type limit**, thus creating a bridge between global algebraic constructions and analytic limit spaces.

Conceptually, this question entwines several profound domains—**Kähler geometry, degeneration theory, Ricci‑flat analysis, Sasaki–Einstein geometry, and metric convergence**—but with an elementary formulation accessible to any specialist in complex differential geometry. Its resolution would yield an explicit algebro‑analytic mechanism realizing asymptotically conical Calabi–Yau metrics as limits of compact ones, illuminating the transition between Fano and Calabi–Yau geometries and potentially advancing understanding of metric degeneration, mirror symmetry, and moduli compactifications.

The problem’s strength lies in the simplicity of its statement versus the depth of the analytic and geometric mechanisms it engages—qualities that make it a quintessentially *good* research-level mathematical problem.